%!TEX program = lualatex
% SPDX-License-Identifier: CC-BY-SA-4.0
% Requires: https://github.com/jemmybutton/fiziko (GPL-3.0)
% Set MPINPUTS to the directory containing fiziko.mp before running lualatex.
\documentclass[11pt]{article}

\usepackage{iftex}
\ifLuaTeX\else
  \PackageError{vintage-latex}{This example requires LuaLaTeX}
    {Run lualatex, not pdflatex or xelatex.}
\fi

\usepackage{fontspec}
\usepackage{unicode-math}
\usepackage{xcolor}
\usepackage{geometry}
\usepackage[tracking=true,protrusion=false,expansion=false]{microtype}
\usepackage{luamplib}

\setmainfont{EB Garamond}[Numbers={OldStyle,Proportional}]
\setmathfont{Garamond-Math.otf}
\definecolor{ragpaper}{HTML}{F3E9D6}
\definecolor{printerink}{HTML}{25201B}
\pagecolor{ragpaper}
\color{printerink}
\geometry{a5paper,top=20mm,bottom=22mm,left=20mm,right=20mm}
\pagestyle{empty}
\setlength{\parindent}{1.1em}
\setlength{\emergencystretch}{.7em}

% fiziko is drawn by LuaTeX's built-in MetaPost. Labels written between
% btex and etex are typeset by this document, so they share its text and math
% fonts; drawoptions tints every stroke with the page's ink.
\mplibforcehmode
\everymplib{
  input fiziko.mp;
  randomseed := 424;
  defineMinStrokeWidth(0.22pt);
  defineLightDirection(-1/8pi, 1/8pi);
  color inkcolor, papercolor;
  inkcolor := (0.145, 0.125, 0.106);
  papercolor := (0.953, 0.914, 0.839);
  drawoptions(withcolor inkcolor);
  % A label on a small patch of paper, so lettering stays legible where it
  % must cross wood grain or hatching.
  vardef clearlabel@#(expr t, z) =
    save p; picture p; p := thelabel@#(t, z);
    fill bbox p withcolor papercolor;
    draw p;
  enddef;
  beginfig(1);
}
\everyendmplib{endfig;}

\begin{document}

\begin{center}
  {\small\textls[120]{\textsc{Proposition VI}}}\par
  \vspace{.3em}
  {\Large Of a Ball upon an Inclined Plane}\par
  \vspace{.2em}
  {\itshape wherein the weight is resolved into its two parts}
\end{center}

\vspace{.9em}

\noindent\textls[70]{\textsc{Let a plane}} be raised at one end until it
makes with the horizon the angle $\theta$, and let a ball of mass $m$ be laid
upon it. The weight $mg$ acts always toward the centre of the earth; but the
plane resists only that part of it which presses directly upon the surface,
and the remainder urges the ball along the slope.

\begin{center}
\begin{mplibcode}
  numeric slopeAngle, planeLength, ballDiameter, ballTravel;
  pair planeFoot, planeTop, ballCenter, contact;
  % Do not call this path "floor": that shadows MetaPost's floor function,
  % which fiziko's wood texture uses.
  path wedge, ground;

  slopeAngle := 24;
  planeLength := 7.4cm;
  ballDiameter := 1.2cm;
  ballTravel := 0.6;

  planeFoot := (0, 0);
  planeTop := planeFoot + planeLength * dir(slopeAngle);

  ground := (planeLength + 1.2cm, 0) -- (-1.2cm, 0);
  draw solidSurface(ground);
  wedge := planeFoot -- (xpart planeTop, 0) -- planeTop -- cycle;
  draw woodenThing(wedge, slopeAngle);

  contact := ballTravel[planeFoot, planeTop];
  ballCenter := contact + 1/2ballDiameter * dir(slopeAngle + 90);
  draw sphere.c(ballDiameter) shifted ballCenter;

  pair gravity, alongSlope, intoSlope;
  gravity := (0, -1.9cm);
  alongSlope := (gravity dotprod dir(slopeAngle)) * dir(slopeAngle);
  intoSlope := gravity - alongSlope;
  drawarrow ballCenter -- ballCenter + gravity withpen pencircle scaled 0.45pt;
  drawarrow ballCenter -- ballCenter + alongSlope
    withpen pencircle scaled 0.45pt;
  drawarrow ballCenter -- ballCenter + intoSlope
    withpen pencircle scaled 0.45pt;
  draw ballCenter + alongSlope -- ballCenter + gravity
    dashed evenly scaled 0.6 withpen pencircle scaled 0.28pt;
  draw ballCenter + intoSlope -- ballCenter + gravity
    dashed evenly scaled 0.6 withpen pencircle scaled 0.28pt;
  clearlabel.rt(btex $mg$ etex, ballCenter + gravity shifted (1mm, 0.6mm));
  clearlabel.ulft(btex $mg\sin\theta$ etex,
    ballCenter + 1/2alongSlope shifted (-0.6mm, 0.6mm));
  clearlabel.lrt(btex $mg\cos\theta$ etex,
    ballCenter + 1/2intoSlope shifted (0.6mm, 0));

  markAngle((xpart planeTop, 0), planeFoot, planeTop)(btex $\theta$ etex);

  draw (xpart planeTop + 0.3cm, 0) -- (xpart planeTop + 0.85cm, 0)
    withpen pencircle scaled 0.3pt;
  draw (xpart planeTop + 0.3cm, ypart planeTop)
    -- (xpart planeTop + 0.85cm, ypart planeTop)
    withpen pencircle scaled 0.3pt;
  drawdblarrow (xpart planeTop + 0.65cm, 0)
    -- (xpart planeTop + 0.65cm, ypart planeTop)
    withpen pencircle scaled 0.4pt;
  clearlabel.rt(btex $h$ etex, (xpart planeTop + 0.75cm, 1/2ypart planeTop));
\end{mplibcode}
\end{center}

\vspace{.4em}
\noindent
The figure shows the ball at the point where its weight is thus divided. The
part $mg\sin\theta$ lies along the plane and alone produces motion; the part
$mg\cos\theta$ is wholly sustained by the wood. If the ball rolls without
slipping, a share of the moving force is spent in turning it, and the
descent is found to proceed with the acceleration
\[
  a = \tfrac{5}{7}\, g \sin\theta ,
\]
which is less than that of a body sliding freely upon a plane without
friction, though independent of the size and the weight of the ball.

\vspace{.4em}
\noindent\footnotesize
The wooden wedge, the hatched ground, the line-shaded sphere, and the angle
mark are fiziko macros drawn by LuaTeX's own MetaPost; the arrows and their
labels are ordinary MetaPost, and the labels are set in the page's fonts.

\vfill
\begin{center}
  {\scriptsize\textls[110]{\textsc{Elements of Mechanics, Book I}}}
\end{center}

\end{document}
